%% kapa: exam=1 round=31 year=2020 period=1 form=A subject=real_estate subject_ko=부동산학원론 seq=11 number=91
%% source: Q-Net 감정평가사 기출문제 (cst00309) artlSeq=5207219 file=extracted/1008/5207219/A형/A형/2020년도 제31회 감정평가사 1차 시험_1교시_A형.pdf
%% conversion: kordoc 4.24.51 -> tools/gen1.py (auto); review=proofread-sonnet; flags=none
\begin{question}{91}
A도시와 B도시 사이에 C도시가 있다. 레일리의 소매인력법칙을 적용할 경우, C도시에서 A도시, B도시로 구매 활동에 유입되는 비율은? (단, C도시의 인구는 모두 A도시 또는 B도시에서 구매하고, 주어진 조건에 한함)
\begin{framedtext}
\begin{tabular}{@{}p{0.45\linewidth}p{0.45\linewidth}@{}}
○ A도시 인구수: 45,000명 & ○ B도시 인구수: 20,000명 \\
○ C도시에서 A도시 간의 거리: 36km & ○ C도시에서 B도시 간의 거리: 18km
\end{tabular}
\end{framedtext}
\begin{choices}
  \item A: 36 \%, B: 64 \%
  \item A: 38 \%, B: 62 \%
  \item A: 40 \%, B: 60 \%
  \item A: 42 \%, B: 58 \%
  \item A: 44 \%, B: 56 \%
\end{choices}
\end{question}
\answer{1}
